\section{One name for one thing}
\label{sec:ch05-one-name-for-one-thing}
\index{Relay!global object identification|(}

Mosaic has had \code{[ID]} on four fields since
chapter~\ref{ch:hotchocolate-quickly}. I put it there because it is what you
write, and I did not check what it did. It did almost nothing, and finding that
out is the most useful hour this chapter cost me.

Here is tag \code{ch04}, the state the last chapter ended in:

\begin{minted}[fontsize=\footnotesize]{graphql}
{ productById(id: "a0000000-0000-4000-8000-000000000001") { id sku } }
\end{minted}

\begin{minted}[fontsize=\footnotesize]{json}
{"data":{"productById":{"id":"a0000000-0000-4000-8000-000000000001",
          "sku":"MOS-FRN-0001"}}}
\end{minted}

The \code{id} field answered with the raw \code{Guid}, and a raw \code{Guid} was
accepted as an \code{[ID]} argument. Whatever the attribute was doing, it was
not encoding anything.

\index{Relay!AddGlobalObjectIdentification}%
One line in \code{Program.cs} changes that, and changes rather more than I
expected:

\begin{minted}{csharp}
builder.AddGraphQL()
    .AddMosaic()
    .AddGlobalObjectIdentification()
\end{minted}

\begin{minted}[fontsize=\footnotesize]{json}
{"data":{"products":[{"id":"UHJvZHVjdDoAAACgAAAAQIAAAAAAAAAB"}, ...]}}
\end{minted}

\index{Relay!opaque identifiers}%
Base64-decode that and you get twenty-four bytes: the ASCII text
\code{Product:} followed by the sixteen raw bytes of the \code{Guid}. The type
name is inside the identifier. That is the whole idea, and it solves a problem
that only exists once a client caches anything: two types can both have a row
numbered one, and a cache keyed on \code{id} alone will confuse them. An
identifier that carries its type cannot collide with an identifier of another
type in the same schema~\autocite{chillicream2026relay}. Hold on to that
qualification. Two services can each mint an identifier for their own
\code{Product}, and what happens when a router has to tell them apart is
chapter~\ref{ch:hard-modeling-problems}'s problem.

\subsection*{The part that is a breaking change}

\index{schema!silent breaking change}%
Send the old query to the new service:

\begin{minted}[fontsize=\footnotesize]{json}
{"errors":[{"message":"The node ID string has an invalid format.",
  "path":["productById"],
  "extensions":{"originalValue":"a0000000-0000-4000-8000-000000000001"}}],
 "data":{"productById":null}}
\end{minted}

Every identifier the service has ever handed out is now unusable, and every
identifier it accepts has a format it did not have yesterday. A client that
stored a product id last week is holding a string that no longer resolves.

Now compare the two schemas. Before: \code{id: ID!}. After: \code{id: ID!}.

\index{schema!what SDL cannot say}%
That is the part worth carrying out of this section. The SDL does not change by
one character, so every tool that compares schemas for breaking changes,
including the one chapter~\ref{ch:ci-cd-for-a-federated-graph} puts in front of
a pull request, reports that nothing happened. The type system has a name for
the shape of an identifier and no name at all for its format, so a change of
format is invisible to it. I would rather have found this by reading than by
sending a query, and I did not.

The lesson generalises past this one directive. Anything a schema treats as an
opaque scalar can change meaning without changing type, and cursors are the
other obvious member of that family. A serialisation you can alter without the
schema moving is a compatibility surface nothing is watching.

So: worth doing anyway, and worth doing early. Mosaic has no external clients,
which is the only condition under which this is a free change, and that
condition expires exactly once.

\subsection*{The interface, and who is allowed in it}

\index{abstract types!Node interface}%
Along with the encoding, that one line adds Mosaic's first abstract type and two
root fields~\autocite{chillicream2026relay}:

\begin{graphqlsdl}
"The node interface is implemented by entities that have a global unique identifier."
interface Node {
  id: ID!
}

type Query {
  node(id: ID!): Node
  nodes(ids: [ID!]!): [Node]!
}
\end{graphqlsdl}

An interface is a promise that several object types share a field, and that a
client holding the interface can select that field without knowing which type it
has~\autocite{graphql2025spec}. \code{Node} is the smallest useful one: a single
field, and the guarantee that anything implementing it can be fetched from that
field alone.

Four of Mosaic's types implement it: \code{Product}, \code{Customer},
\code{Review} and \code{Order}. One deliberately does not, and the exclusion
teaches more than the inclusions.

\index{abstract types!what does not belong}%
\code{OrderLine} has no identifier. Not a hidden one, not an inconvenient one:
the domain does not have a concept of a line's identity apart from the order it
sits in, which is why chapter~\ref{ch:data-without-the-n-plus-1} gave the table
a shadow key rather than inventing a property nobody would use. Putting
\code{OrderLine} in \code{Node} would mean promising a client that a line can be
fetched on its own, and the honest answer is that there is nothing to fetch it
by. The interface is a claim about the domain and not a badge to hand out.

\subsection*{The bill the interface presents}

\index{Relay!refetch cost}%
Making four types refetchable had a consequence I did not anticipate and should
have. Two of the four had no way to be fetched by their own identifier. The only
route to an \code{Order} was through the customer who placed it; the only route
to a \code{Review} was through the product it was written about. Both needed a
new service method and a new DataLoader:

\begin{minted}[fontsize=\footnotesize]{csharp}
public async Task<IReadOnlyDictionary<Guid, Order>> GetOrdersByIdsAsync(
    IReadOnlyList<Guid> ids,
    CancellationToken cancellationToken)
{
    counter.RecordLookup();

    return await db.Orders
        .AsNoTracking()
        .Where(o => ids.Contains(o.Id))
        .ToDictionaryAsync(o => o.Id, cancellationToken);
}
\end{minted}

A batch rather than a single lookup, because \code{nodes(ids:)} can ask for many
at once and resolving them one at a time would rebuild the problem the last
chapter spent itself removing.

That bill is worth naming, because it is the same bill twice. The Node interface
says: give me an identifier and I will give you the thing. From
chapter~\ref{ch:the-first-cut-extracting-catalog} onward, federation says
exactly the same, in almost the same words, to a service that no longer holds
the data the identifier refers to. A codebase that has already answered the
question for its own graph has done most of the work of answering it for a
router.

\subsection*{Two routes to a node resolver, and the documentation knows one}

\index{HotChocolate!NodeResolver}%
The documentation describes a \code{[Node]} attribute on the class, with the
resolver found by naming convention: a static method called \code{Get} or
\code{GetAsync}, or either of those with the type name after the
\code{Get}~\autocite{chillicream2026relay}. That route exists and is
implemented by reflection in \code{NodeAttribute}.

There is a second route, it is the one the source generator implements, and it
appears nowhere in the documentation. Mark a method \code{[NodeResolver]}
inside a class that already carries \code{[ObjectType<T>]}, and
\code{ObjectTypeFileBuilder} emits this into the generated configuration:

\begin{minted}[fontsize=\footnotesize]{csharp}
descriptor
    .ImplementsNode()
    .ResolveNode(resolvers.ResolveProductAsync().Resolver!);
\end{minted}

No \code{[Node]} attribute, no naming convention. Mosaic uses it because every
domain type is already an \code{[ObjectType<T>]} partial class, so the node
resolver lives beside the other fields of that type, in the folder of the domain
that owns it:

\begin{minted}[fontsize=\footnotesize]{csharp}
[NodeResolver]
public static async Task<Product?> ResolveProductAsync(
    Guid id,
    IProductByIdDataLoader productById,
    CancellationToken cancellationToken)
    => await productById.LoadAsync(id, cancellationToken);
\end{minted}

\index{HotChocolate!analyzer diagnostics}%
The generator enforces four rules that the documentation does not state, and
each has its own diagnostic. The first parameter must be named \code{id},
exactly, or you get HC0104. It must not carry \code{[ID]}, because a node
resolver already declares it as one, or HC0092. The method must be public, or
HC0093. And \code{id} must be the only field argument, or HC0083.

I list them because every one of them is a compile-time error, which is the
whole reason this route is the better one. The reflection route resolves its
convention when the schema is built, so a method named wrongly is a runtime
surprise. Here the compiler tells you before the service starts. That is worth
more than the attribute being shorter.

\subsection*{Refetching, end to end}

\index{Relay!node field}%
Which leaves the field the whole arrangement exists for. Hand it one of those
identifiers and nothing else:

\begin{minted}[fontsize=\footnotesize]{graphql}
{
  node(id: "UHJvZHVjdDoAAACgAAAAQIAAAAAAAAAB") {
    __typename
    ... on Product { title sku }
  }
}
\end{minted}

\begin{minted}[fontsize=\footnotesize]{json}
{"data":{"node":{"__typename":"Product","title":"Larsen Oak Dining Table",
          "sku":"MOS-FRN-0001"}}}
\end{minted}

The inline fragment is not optional. \code{node} returns \code{Node}, and a
\code{Node} has exactly one field, so anything beyond \code{id} has to be asked
for on a concrete type. That is the cost of an abstract type and it is the same
cost every time: the client does the narrowing, in the query, where it is
visible.

\code{\_\_typename} is doing real work in that document rather than decorating it.
It is how a client that asked \code{node} for something knows what came back,
and it is the field every cache in the GraphQL ecosystem keys on alongside
\code{id}. The two together are why the encoding at the top of this section
matters.

Which is one field made honest. The next one changes shape entirely.
\index{Relay!global object identification|)}
